% !TEX root = ../main.tex
%% Chapitre 13 — Matrices
%% Programme officiel MPSI/MP2I (BO spécial n°1 du 11/02/2021), deuxième semestre,
%% section « Matrices » : A) matrices et applications linéaires (matrice d'une AL,
%% application canoniquement associée, systèmes linéaires) ; B) changements de
%% bases, matrices équivalentes et rang (J_r, opérations élémentaires),
%% matrices semblables et trace.
%% Résultats vérifiés par verifications/chapitre_13.py.

\begin{ouverture}[passage]{13}{Matrices}{Partie 2 · Deuxième semestre}
  {Matrice d'une application linéaire\enspace·\enspace Rang\enspace·\enspace
   Changements de bases\enspace·\enspace Équivalence\enspace·\enspace Similitude\enspace·\enspace Trace}
  \objectif{Représenter vecteurs, familles et applications linéaires par des matrices.}
  \objectif{Traduire noyau, image, rang et systèmes linéaires dans le langage matriciel.}
  \objectif{Manier les matrices de passage et les formules de changement de base.}
  \objectif{Caractériser le rang par l'équivalence à $J_r$ et le calculer par opérations élémentaires.}
  \objectif{Reconnaître des matrices semblables ; utiliser la trace comme invariant.}
  \objectif{Chercher une base dans laquelle un endomorphisme devient simple.}
  \prerequis{Calcul matriciel et systèmes linéaires : produit, inverse, pivot (chapitre 9).}
  \prerequis{Espaces vectoriels : bases, dimension, applications linéaires, rang (chapitre 12).}
  \prerequis{Nombres complexes (chapitre 3) pour l'exemple de $\Cx$ vu comme plan réel.}
\end{ouverture}

%% =====================================================================
%%  L'ESSENTIEL
%% =====================================================================
\essentiel

\essSection{1}{La matrice d'une application linéaire}

\begin{definition}[Matrices d'un vecteur, d'une famille, d'une application]
Soit $e=(e_1,\dots,e_p)$ une base de $E$ et $f=(f_1,\dots,f_n)$ une base de
$F$. La \textbf{matrice d'un vecteur} $x$ dans $e$ est la colonne de ses
coordonnées ; celle d'une famille, la juxtaposition des colonnes. La
\textbf{matrice de} $u\in\mathcal{L}(E,F)$ \textbf{dans les bases} $e,f$ est
\[
\operatorname{Mat}_{e,f}(u)=\bigl(\text{colonnes des }u(e_j)\text{ dans }f\bigr)
\in\mathcal{M}_{n,p}(\K) :
\]
la $j$-ième colonne contient les coordonnées de $u(e_j)$. Pour un
endomorphisme, on prend la même base au départ et à l'arrivée :
$\operatorname{Mat}_{e}(u)$.
\end{definition}

\begin{theoreme}[Le dictionnaire complet]
L'application $u\mapsto\operatorname{Mat}_{e,f}(u)$ est un
\emph{isomorphisme} de $\mathcal{L}(E,F)$ sur $\mathcal{M}_{n,p}(\K)$, et :
\begin{itemize}[label={\color{otNavy}\textbullet},leftmargin=1.4em,itemsep=0.6mm,topsep=0.8mm]
  \item si $X$ est la colonne de $x$ et $Y$ celle de $u(x)$, alors
    $Y=AX$ avec $A=\operatorname{Mat}_{e,f}(u)$ ;
  \item la composée devient le produit :
    $\operatorname{Mat}(v\circ u)=\operatorname{Mat}(v)\operatorname{Mat}(u)$
    (bases intermédiaires cohérentes) — pour les endomorphismes,
    $u\mapsto\operatorname{Mat}_e(u)$ est un isomorphisme d'\emph{anneaux}
    de $\mathcal{L}(E)$ sur $\mathcal{M}_n(\K)$ ;
  \item $u$ est un isomorphisme si et seulement si sa matrice est
    inversible, et $\operatorname{Mat}(u^{-1})=\operatorname{Mat}(u)^{-1}$.
\end{itemize}
\end{theoreme}

\begin{exemple}[$\Cx$ vu comme plan réel]
Dans la base $(1,\mathrm{i})$ de $\Cx$ sur $\R$, la multiplication par
$a+\mathrm{i}b$ a pour matrice
$\begin{pmatrix}a&-b\\b&a\end{pmatrix}$ — on retrouve le sous-anneau du
chapitre 9. La multiplication par $\mathrm{i}$ devient la rotation
$J=\begin{pmatrix}0&-1\\1&0\end{pmatrix}$, avec $J^{2}=-I_2$ : l'identité
$\mathrm{i}^{2}=-1$ devient une égalité de matrices, et pour
$a+\mathrm{i}b=\cos\theta+\mathrm{i}\sin\theta$ on obtient la matrice de la
rotation d'angle $\theta$.
\end{exemple}

\essSection{2}{La vision matricielle : noyau, image, rang, systèmes}

\begin{theoreme}[Application linéaire canoniquement associée]
À toute $A\in\mathcal{M}_{n,p}(\K)$ on associe $X\mapsto AX$ de $\K^{p}$
dans $\K^{n}$ (en identifiant $\mathcal{M}_{n,1}$ et $\K^{n}$). On définit
ainsi $\operatorname{Ker}A$, $\operatorname{Im}A$ et
$\operatorname{rg}A$ : \emph{les colonnes engendrent l'image, les lignes
donnent un système d'équations du noyau}. Pour $A$ carrée d'ordre $n$,
sont équivalents :
\[
A\ \text{inversible}
\iff\operatorname{Ker}A=\ens{0}
\iff\text{les colonnes engendrent }\K^{n}
\iff\operatorname{rg}A=n ,
\]
et toute matrice carrée inversible à gauche \emph{ou} à droite est
inversible. On retrouve ainsi qu'une triangulaire est inversible si et
seulement si sa diagonale ne s'annule pas.
\end{theoreme}

\begin{propriete}[Systèmes linéaires]
Le système homogène $AX=0$ a pour ensemble de solutions le sous-espace
$\operatorname{Ker}A$, de dimension $p-\operatorname{rg}A$ (théorème du
rang). Le système $AX=B$ est compatible si et seulement si
$B\in\operatorname{Im}A$, et ses solutions forment alors le sous-espace
affine $X_0+\operatorname{Ker}A$. Si $A$ est carrée inversible, le système
est dit \emph{de Cramer} : unique solution $X=A^{-1}B$.
\end{propriete}

\essSection{3}{Changements de bases}

\noindent
\begin{minipage}[c]{0.55\linewidth}
\begin{definition}[Matrice de passage]
Si $e$ et $e'$ sont deux bases de $E$, la \textbf{matrice de passage}
$P=P_{e\to e'}$ est la matrice de la famille $e'$ dans la base $e$
(colonnes = nouveaux vecteurs lus dans l'ancienne base). Elle est
inversible, et $P^{-1}=P_{e'\to e}$.
\end{definition}
\begin{theoreme}[Formules de changement de base]
Avec $X=\operatorname{Mat}_e(x)$ et $X'=\operatorname{Mat}_{e'}(x)$ :
\[
X=PX' .
\]
Pour $u\in\mathcal{L}(E,F)$, $A=\operatorname{Mat}_{e,f}(u)$ et
$B=\operatorname{Mat}_{e',f'}(u)$ : $B=Q^{-1}AP$ avec $P=P_{e\to e'}$,
$Q=P_{f\to f'}$. Pour un endomorphisme :
\[
B=P^{-1}AP .
\]
\end{theoreme}
\end{minipage}\hfill
\begin{minipage}[c]{0.41\linewidth}
\centering
\begin{tikzpicture}[scale=1.02,line cap=round]
  \draw[otMuted,line width=0.8pt,-{Stealth[length=2mm]}] (0,0) -- (2.5,0);
  \draw[otMuted,line width=0.8pt,-{Stealth[length=2mm]}] (0,0) -- (0,2.5);
  \draw[otGreenDark,line width=0.8pt,-{Stealth[length=2mm]}] (0,0) -- (2.21,1.17);
  \draw[otGreenDark,line width=0.8pt,-{Stealth[length=2mm]}] (0,0) -- (-1.17,2.21);
  \draw[otMuted,line width=0.5pt,dash pattern=on 2pt off 1.5pt]
    (1.5,1.9) -- (1.5,0) (1.5,1.9) -- (0,1.9);
  \draw[otGreenDark,opacity=0.8,line width=0.5pt,dash pattern=on 2pt off 1.5pt]
    (1.5,1.9) -- (1.94,1.03) (1.5,1.9) -- (-0.44,0.83);
  \draw[otInk,line width=1.1pt,-{Stealth[length=2.2mm]}] (0,0) -- (1.5,1.9);
  \node[font=\scriptsize,text=otInk,anchor=south] at (1.5,1.95) {$x$};
  \node[font=\scriptsize,text=otMuted] at (2.5,-0.25) {$e_1$};
  \node[font=\scriptsize,text=otMuted] at (-0.25,2.45) {$e_2$};
  \node[font=\scriptsize,text=otGreenDark] at (2.5,1.25) {$f_1$};
  \node[font=\scriptsize,text=otGreenDark] at (-1.35,2.4) {$f_2$};
\end{tikzpicture}\\[0.5mm]
\legende{Un même vecteur, deux jeux de coordonnées : la matrice de passage
fait la traduction $X=PX'$.}
\end{minipage}

\begin{attention}
Dans $X=PX'$, la matrice de passage $P$ (vers la \emph{nouvelle} base)
transforme les \emph{nouvelles} coordonnées en \emph{anciennes} — le sens
inverse de ce qu'on attend naïvement. Un moyen mnémotechnique : les
colonnes de $P$ vivent dans l'ancienne base, donc $PX'$ aussi.
\end{attention}

\begin{savoirfaire}[Rendre un endomorphisme simple]
Pour étudier $u$ de matrice $A$, chercher une base adaptée (vecteurs
invariants, directions privilégiées…) et poser $P$ la matrice de passage :
$B=P^{-1}AP$ y est souvent diagonale ou triangulaire. Exemple : le
projecteur $p(x,y)=\bigl(\frac{x+y}{2},\frac{x+y}{2}\bigr)$ du chapitre 12
a pour matrice $\frac12\begin{pmatrix}1&1\\1&1\end{pmatrix}$ dans la base
canonique, mais $\begin{pmatrix}1&0\\0&0\end{pmatrix}$ dans la base
$\bigl((1,1),(1,-1)\bigr)$ adaptée à
$\operatorname{Im}p\oplus\operatorname{Ker}p$. Les puissances $A^{n}$ se
calculent alors en une ligne (exercice 9).
\end{savoirfaire}

\essSection{4}{Matrices équivalentes et rang}

\begin{theoreme}[La forme normale $J_r$]
Si $u\in\mathcal{L}(E,F)$ est de rang $r$, il existe des bases de $E$ et
$F$ dans lesquelles sa matrice est
\[
J_r=\begin{pmatrix}I_r&0\\0&0\end{pmatrix} .
\]
Deux matrices $A,B\in\mathcal{M}_{n,p}(\K)$ sont dites
\textbf{équivalentes} s'il existe $Q\in\operatorname{GL}_n$,
$P\in\operatorname{GL}_p$ avec $B=QAP$ : cela revient à représenter la même
application dans d'autres bases. \emph{Une matrice est de rang $r$ si et
seulement si elle est équivalente à $J_r$} : le rang classifie entièrement
les matrices à équivalence près. Conséquence remarquable :
$\operatorname{rg}\bigl(A^{\mathsf T}\bigr)=\operatorname{rg}A$.
\end{theoreme}

\begin{propriete}[Rang et matrices extraites ; opérations élémentaires]
Le rang d'une matrice extraite est au plus celui de $A$, et
$\operatorname{rg}A$ est la taille maximale d'une matrice carrée extraite
inversible. Les opérations élémentaires sur les colonnes conservent
l'\emph{image}, celles sur les lignes conservent le \emph{noyau} : toutes
conservent le rang. En pratique, on calcule donc le rang par le pivot de
Gauss (chapitre 9), en comptant les pivots non nuls.
\end{propriete}

\essSection{5}{Matrices semblables et trace}

\begin{definition}[Similitude]
$A,B\in\mathcal{M}_n(\K)$ sont \textbf{semblables} s'il existe
$P\in\operatorname{GL}_n(\K)$ telle que $B=P^{-1}AP$ : ce sont les matrices
d'un \emph{même endomorphisme} dans deux bases. La similitude est plus
fine que l'équivalence : semblables $\Rightarrow$ équivalentes (même rang),
mais la réciproque est fausse.
\end{definition}

\begin{theoreme}[Trace]
La \textbf{trace} $\operatorname{tr}A$ est la somme des coefficients
diagonaux. Elle est linéaire et vérifie l'identité fondamentale
\[
\operatorname{tr}(AB)=\operatorname{tr}(BA),
\qquad\text{d'où}\qquad
\operatorname{tr}\bigl(P^{-1}AP\bigr)=\operatorname{tr}A :
\]
deux matrices semblables ont la même trace. On définit ainsi la trace
$\operatorname{tr}u$ d'un endomorphisme (indépendante de la base), linéaire
avec $\operatorname{tr}(uv)=\operatorname{tr}(vu)$. Pour un projecteur :
$\operatorname{tr}p=\operatorname{rg}p$ — la trace \emph{compte} la
dimension de l'image.
\end{theoreme}

\begin{savoirfaire}[Montrer que deux matrices ne sont pas semblables]
Comparer les invariants de similitude : le rang, la trace (plus tard le
déterminant, chapitre 14). Ainsi
$N=\begin{pmatrix}0&1\\0&0\end{pmatrix}$ et la matrice nulle ont la même
trace mais pas le même rang : non semblables. Attention, ces invariants ne
suffisent pas : $I_2$ et $\begin{pmatrix}1&1\\0&1\end{pmatrix}$ ont même
rang et même trace sans être semblables (exercice 14). Pour montrer que
deux matrices \emph{sont} semblables, exhiber la matrice de passage ou la
base adaptée.
\end{savoirfaire}

%% =====================================================================
%%  EXERCICES
%% =====================================================================
\exercices
\rubrique{Applications}

\begin{exercice}[1]{Écrire des matrices}
Donner la matrice, dans les bases canoniques, de :
\begin{questions}
  \item $u:\R^{2}\to\R^{3}$, $u(x,y)=(x+y,\ 2x-y,\ x)$ ;
  \item la dérivation $D:P\mapsto P'$ sur $\R_2[X]$ ;
  \item $T:P\mapsto XP'$ sur $\R_2[X]$. Que remarque-t-on ?
\end{questions}
\end{exercice}

\begin{exercice}[1]{Produits et composées}
Soient $A=\begin{pmatrix}1&2\\3&4\end{pmatrix}$,
$B=\begin{pmatrix}0&1\\1&0\end{pmatrix}$ et
$C=\begin{pmatrix}1&0\\2&1\end{pmatrix}$.
\begin{questions}
  \item Calculer $AB$ et $BA$. Que constate-t-on ?
  \item Vérifier sur cet exemple que $(AB)C=A(BC)$.
  \item Si $A$ et $B$ sont les matrices de $u$ et $v$ dans la base
    canonique, quelle est la matrice de $u\circ v$ ? de $v\circ u$ ?
\end{questions}
\end{exercice}

\begin{exercice}[1]{$Y=AX$ : images et antécédents}
Soit $A=\begin{pmatrix}1&1\\1&2\end{pmatrix}$, matrice de $u$ dans la base
canonique de $\R^{2}$.
\begin{questions}
  \item Calculer l'image de $(1,2)$.
  \item Déterminer tous les antécédents de $(3,5)$.
  \item Pourquoi chaque vecteur a-t-il exactement un antécédent ?
\end{questions}
\end{exercice}

\begin{exercice}[1]{Noyau, image, rang d'une matrice}
Soit $A=\begin{pmatrix}1&2&3\\2&4&6\\1&1&1\end{pmatrix}$.
\begin{questions}
  \item Déterminer $\operatorname{rg}A$.
  \item Donner une base de $\operatorname{Ker}A$ et une base de
    $\operatorname{Im}A$.
  \item Vérifier le théorème du rang.
\end{questions}
\end{exercice}

\begin{exercice}[1]{Inversible ou pas ?}
\begin{questions}
  \item Montrer que $A=\begin{pmatrix}1&2\\3&4\end{pmatrix}$ est
    inversible et calculer $A^{-1}$.
  \item La matrice $B=\begin{pmatrix}1&2&3\\4&5&6\\7&8&9\end{pmatrix}$
    est-elle inversible ?
  \item À quelle condition une matrice triangulaire est-elle inversible ?
    L'appliquer à $T=\begin{pmatrix}2&1&5\\0&3&-1\\0&0&4\end{pmatrix}$.
\end{questions}
\end{exercice}

\begin{exercice}[1]{Matrice de passage : premiers pas}
Dans $\R^{2}$, on pose $e'_1=(1,1)$ et $e'_2=(1,-1)$.
\begin{questions}
  \item Justifier que $e'=(e'_1,e'_2)$ est une base et écrire la matrice
    de passage $P$ de la base canonique à $e'$.
  \item Calculer $P^{-1}$.
  \item En déduire les coordonnées de $(3,1)$ dans $e'$, et vérifier.
\end{questions}
\end{exercice}

\begin{exercice}[1]{Un endomorphisme dans deux bases}
Soit $p$ le projecteur de matrice
$A=\dfrac12\begin{pmatrix}1&1\\1&1\end{pmatrix}$ dans la base canonique
(chapitre 12).
\begin{questions}
  \item Rappeler pourquoi $\bigl((1,1),(1,-1)\bigr)$ est adaptée à la
    décomposition $\R^{2}=\operatorname{Im}p\oplus\operatorname{Ker}p$.
  \item Calculer $P^{-1}AP$ et commenter.
\end{questions}
\end{exercice}

\begin{exercice}[1]{La trace : premiers calculs}
Avec les matrices $A$ et $B$ de l'exercice 2 :
\begin{questions}
  \item Calculer $\operatorname{tr}A$, $\operatorname{tr}(AB)$ et
    $\operatorname{tr}(BA)$.
  \item La trace est-elle multiplicative
    ($\operatorname{tr}(AB)=\operatorname{tr}A\operatorname{tr}B$) ?
  \item Calculer la trace du projecteur de l'exercice 7 et comparer à son
    rang.
\end{questions}
\end{exercice}

\rubrique{Entraînement}

\begin{exercice}[2]{Puissances par changement de base}
Soit $A=\begin{pmatrix}1&2\\2&1\end{pmatrix}$.
\begin{questions}
  \item Vérifier que $(1,1)$ et $(1,-1)$ sont envoyés sur des multiples
    d'eux-mêmes, et en déduire que $D=P^{-1}AP$ est diagonale pour
    $P=\begin{pmatrix}1&1\\1&-1\end{pmatrix}$.
  \item En déduire, pour tout $n\in\N$ :
    \[
    A^{n}=\frac12\begin{pmatrix}3^{n}+(-1)^{n}&3^{n}-(-1)^{n}\\
    3^{n}-(-1)^{n}&3^{n}+(-1)^{n}\end{pmatrix}.
    \]
\end{questions}
\end{exercice}

\begin{exercice}[2]{Le rang par opérations élémentaires}
\begin{questions}
  \item Calculer le rang de
    $\begin{pmatrix}1&2&3&4\\2&4&6&8\\1&1&1&1\end{pmatrix}$.
  \item Pour $t\in\R$, déterminer le rang de
    $M_t=\begin{pmatrix}1&1&t\\1&t&1\\t&1&1\end{pmatrix}$ suivant les
    valeurs de $t$.
\end{questions}
\end{exercice}

\begin{exercice}[2]{Équivalence et forme normale $J_r$}
Soit $A=\begin{pmatrix}1&2&3\\2&4&6\end{pmatrix}$.
\begin{questions}
  \item Déterminer $\operatorname{rg}A$, et expliciter
    $Q\in\operatorname{GL}_2(\R)$ et $P\in\operatorname{GL}_3(\R)$ telles
    que $QAP=J_1$.
  \item Pourquoi deux matrices de $\mathcal{M}_{3}(\R)$ de même rang
    sont-elles toujours équivalentes ? Sont-elles toujours semblables ?
  \item En déduire $\operatorname{rg}\bigl(A^{\mathsf T}\bigr)$.
\end{questions}
\end{exercice}

\begin{exercice}[2]{Systèmes de Cramer}
\begin{questions}
  \item Résoudre $\begin{cases}2x+y=5\\x+3y=5\end{cases}$ en calculant
    l'inverse de la matrice du système.
  \item Montrer que pour tout $(b_1,b_2,b_3)$, le système
    $x+y+z=b_1$, $x+2y+4z=b_2$, $x+3y+9z=b_3$ possède une unique
    solution. Quel lien avec l'interpolation (chapitre 12) ?
\end{questions}
\end{exercice}

\begin{exercice}[2]{$\Cx$ vu comme plan réel, le retour}
Pour $(a,b)\in\R^{2}$, soit $M(a,b)=\begin{pmatrix}a&-b\\b&a\end{pmatrix}$
et $J=M(0,1)$.
\begin{questions}
  \item Interpréter $M(a,b)$ comme la matrice de $z\mapsto(a+\mathrm{i}b)z$
    dans la base $(1,\mathrm{i})$, et retrouver $J^{2}=-I_2$.
  \item Montrer que $\det M(a,b)=a^{2}+b^{2}=|a+\mathrm{i}b|^{2}$.
  \item Pour $a+\mathrm{i}b=r(\cos\theta+\mathrm{i}\sin\theta)$,
    reconnaître $M(a,b)$ comme une similitude directe : rotation d'angle
    $\theta$ composée avec l'homothétie de rapport $r$.
\end{questions}
\end{exercice}

\begin{exercice}[2]{Semblables ou pas ?}
\begin{questions}
  \item Montrer que $N=\begin{pmatrix}0&1\\0&0\end{pmatrix}$ et
    $N'=\begin{pmatrix}0&0\\1&0\end{pmatrix}$ sont semblables.
  \item Montrer que $I_2$ n'est semblable qu'à elle-même. La matrice
    $\begin{pmatrix}1&1\\0&1\end{pmatrix}$ est-elle semblable à $I_2$ ?
  \item $\begin{pmatrix}1&0\\0&0\end{pmatrix}$ et $I_2$ sont-elles
    semblables ? Et $N$ et la matrice nulle ? Quels invariants utilise-t-on ?
\end{questions}
\end{exercice}

\begin{exercice}[2]{La dérivation et la base de Taylor}
Soit $D:P\mapsto P'$ sur $\R_3[X]$.
\begin{questions}
  \item Écrire la matrice $A$ de $D$ dans la base canonique
    $(1,X,X^{2},X^{3})$, et donner son rang et son noyau.
  \item Écrire la matrice de $D$ dans la base de Taylor
    $\bigl(1,X-1,(X-1)^{2},(X-1)^{3}\bigr)$. Qu'observe-t-on, et
    pourquoi était-ce prévisible ?
\end{questions}
\end{exercice}

\begin{exercice}[2]{L'identité $\operatorname{tr}(AB)=\operatorname{tr}(BA)$}
\begin{questions}
  \item Démontrer que $\operatorname{tr}(AB)=\operatorname{tr}(BA)$ pour
    toutes $A\in\mathcal{M}_{n,p}(\K)$, $B\in\mathcal{M}_{p,n}(\K)$
    (écrire les deux doubles sommes).
  \item En déduire que deux matrices semblables ont la même trace, et que
    la trace d'un endomorphisme est bien définie.
  \item Montrer que la trace d'un projecteur vaut son rang.
  \item Existe-t-il $P$ inversible telle que $P^{-1}AP=A+I_n$ ?
\end{questions}
\end{exercice}

\begin{exercice}[2]{Rang et matrices extraites}
Soit $A=\begin{pmatrix}1&2&3&4\\2&4&6&8\\1&1&1&1\end{pmatrix}$ (exercice 10).
\begin{questions}
  \item Exhiber une matrice extraite $2\times2$ inversible. Que peut-on en
    déduire sur $\operatorname{rg}A$ ?
  \item Vérifier que les quatre matrices extraites $3\times3$ sont toutes
    non inversibles, et conclure.
\end{questions}
\end{exercice}

\begin{exercice}[2]{Matrices de rang 1}
Soit $A\in\mathcal{M}_n(\K)$ de rang $1$.
\begin{questions}
  \item Montrer qu'il existe des colonnes non nulles $U,V\in\K^{n}$ telles
    que $A=UV^{\mathsf T}$.
  \item En déduire l'identité $A^{2}=\operatorname{tr}(A)\,A$, et la
    vérifier sur $A=\begin{pmatrix}2&1&1\\4&2&2\\6&3&3\end{pmatrix}$.
  \item À quelle condition une matrice de rang $1$ est-elle nilpotente ?
\end{questions}
\end{exercice}

\rubrique{Approfondissement}

\begin{exercice}[3]{Le commutant d'une matrice diagonale}
Soit $D=\operatorname{diag}(1,2,3)$ et
$\mathcal{C}=\ens{M\in\mathcal{M}_3(\R)\mid MD=DM}$.
\begin{questions}
  \item Montrer que $\mathcal{C}$ est un sous-espace vectoriel de
    $\mathcal{M}_3(\R)$.
  \item Montrer que $\mathcal{C}$ est exactement l'ensemble des matrices
    diagonales, et donner sa dimension.
  \item Que se passerait-il si deux coefficients de $D$ étaient égaux ?
\end{questions}
\end{exercice}

\begin{exercice}[3]{Le centre de $\mathcal{M}_n(\K)$}
Soit $A\in\mathcal{M}_n(\K)$ commutant avec \emph{toutes} les matrices.
\begin{questions}
  \item En faisant commuter $A$ avec les matrices élémentaires $E_{ij}$,
    montrer que $A$ est scalaire : $A=\lambda I_n$.
  \item En déduire le centre de l'anneau $\mathcal{M}_n(\K)$, puis les
    endomorphismes qui commutent avec tous les autres.
\end{questions}
\end{exercice}

\begin{exercice}[3]{Une base où tout devient simple}
Soit $A=\begin{pmatrix}1&1\\-1&3\end{pmatrix}$.
\begin{questions}
  \item Vérifier que $(A-2I_2)^{2}=0$ ; en déduire que $2$ est la seule
    valeur possible pour $\lambda$ tel que $A-\lambda I_2$ soit non
    inversible.
  \item Trouver $v$ tel que $Av=2v$, puis $w$ tel que $(A-2I_2)w=v$, et
    montrer que $P^{-1}AP=\begin{pmatrix}2&1\\0&2\end{pmatrix}$ pour
    $P=\operatorname{Mat}(v,w)$.
  \item En écrivant $A=2I_2+N$ avec $N^{2}=0$, calculer $A^{n}$ pour tout
    $n\in\N$.
\end{questions}
\end{exercice}

\begin{exercice}[3]{$\operatorname{tr}\bigl(AA^{\mathsf T}\bigr)$, un produit scalaire avant l'heure}
Pour $A,B\in\mathcal{M}_n(\R)$, on pose
$\langle A,B\rangle=\operatorname{tr}\bigl(AB^{\mathsf T}\bigr)$.
\begin{questions}
  \item Montrer que
    $\langle A,A\rangle=\sum_{i,j}a_{ij}^{2}$, et en déduire que
    $\langle A,A\rangle=0$ si et seulement si $A=0$.
  \item Vérifier la symétrie $\langle A,B\rangle=\langle B,A\rangle$ et la
    bilinéarité.
  \item Montrer l'inégalité
    $\langle A,B\rangle^{2}\leqslant\langle A,A\rangle\langle B,B\rangle$
    en étudiant le trinôme $t\mapsto\langle A+tB,A+tB\rangle$ — on y
    reviendra au chapitre 18.
\end{questions}
\end{exercice}

\begin{exercice}[3]{Le théorème $J_r$, démonstration guidée}
Soit $u\in\mathcal{L}(E,F)$ de rang $r$, avec $\dim E=p$, $\dim F=n$.
\begin{questions}
  \item Justifier l'existence d'un supplémentaire $S$ de
    $\operatorname{Ker}u$ dans $E$, avec $\dim S=r$, et montrer que $u$
    induit un isomorphisme de $S$ sur $\operatorname{Im}u$.
  \item Construire une base $(s_1,\dots,s_r,k_1,\dots,k_{p-r})$ de $E$ et
    une base de $F$ commençant par $\bigl(u(s_1),\dots,u(s_r)\bigr)$, et
    vérifier que la matrice de $u$ dans ces bases est $J_r$.
  \item En déduire que toute matrice de rang $r$ est équivalente à $J_r$,
    puis que $\operatorname{rg}\bigl(A^{\mathsf T}\bigr)=\operatorname{rg}A$.
\end{questions}
\end{exercice}

\begin{exercice}[3]{Fibonacci en matrice}
Soit $F=\begin{pmatrix}1&1\\1&0\end{pmatrix}$ et $(F_n)$ la suite de
Fibonacci ($F_0=0$, $F_1=1$).
\begin{questions}
  \item Montrer par récurrence que
    $F^{n}=\begin{pmatrix}F_{n+1}&F_n\\F_n&F_{n-1}\end{pmatrix}$ pour
    $n\geqslant1$.
  \item En prenant le déterminant, retrouver l'identité de Cassini
    $F_{n+1}F_{n-1}-F_n^{2}=(-1)^{n}$ (chapitre 5), et en déduire que deux
    termes consécutifs sont premiers entre eux (chapitre 7).
  \item Les directions $(\varphi,1)$ et $(\psi,1)$, où $\varphi,\psi$ sont
    les racines de $x^{2}=x+1$, sont stables par $F$ : faire le lien avec
    la base de suites géométriques du chapitre 12 et la formule de Binet.
\end{questions}
\end{exercice}

% !TEX root = ../../main.tex
%% Défi du chapitre 13 — Les matrices 2×2 à similitude près
%% Résultats vérifiés par verifications/defis_11_15.py.

\rubrique{Défi}

\begin{defi}{Les matrices $2\times2$ à similitude près}{1 h 30}
On travaille dans $\mathcal{M}_2(\K)$ avec $\K=\R$ ou $\Cx$. Rang et
trace sont des invariants de similitude, le déterminant aussi
(chapitre 14). On va montrer qu'en taille $2$, \emph{trace et déterminant
suffisent} à décider si deux matrices non scalaires sont semblables.

\partieExo{A}{La forme compagnon}
\begin{questions}
  \item Montrer que pour tout $A\in\mathcal{M}_2(\K)$ :
    $A^{2}-\operatorname{tr}(A)\,A+\det(A)\,I_2=0$.
  \item On suppose que pour tout vecteur non nul $x\in\K^{2}$, la famille
    $(x,Ax)$ est liée. En appliquant cette hypothèse à $e_1$, $e_2$ et
    $e_1+e_2$, montrer que $A$ est \emph{scalaire} : $A=\lambda I_2$.
  \item On suppose $A$ non scalaire. Montrer qu'il existe $x$ tel que
    $(x,Ax)$ soit une base de $\K^{2}$, et que dans cette base
    l'endomorphisme canoniquement associé à $A$ a pour matrice
    \[
    C_A=\begin{pmatrix}0&-\det A\\ 1&\operatorname{tr}A\end{pmatrix}
    \qquad(\text{matrice \emph{compagnon}}).
    \]
\end{questions}

\partieExo{B}{La classification}
\begin{questions}[resume]
  \item En déduire que deux matrices \emph{non scalaires} de
    $\mathcal{M}_2(\K)$ sont semblables si et seulement si elles ont même
    trace et même déterminant. Qu'en est-il d'une matrice scalaire ?
  \item Étudier si les matrices suivantes sont semblables, et donner une
    matrice de passage quand c'est le cas :
    \begin{itemize}[label={\color{otNavy}\textbullet},leftmargin=1.2em,itemsep=0.5mm,topsep=0.6mm]
      \item $T=\begin{pmatrix}1&1\\0&1\end{pmatrix}$ et
        $\begin{pmatrix}2&1\\-1&0\end{pmatrix}$ ;
      \item $R=\begin{pmatrix}0&-1\\1&0\end{pmatrix}$ et
        $S=\begin{pmatrix}1&-2\\1&-1\end{pmatrix}$ ;
      \item $T$ et $I_2$.
    \end{itemize}
\end{questions}

\partieExo{C}{Puissances et suites récurrentes}
\begin{questions}[resume]
  \item Soit $A$ non scalaire avec $\operatorname{tr}A=2$ et $\det A=1$.
    Montrer que $(A-I_2)^{2}=0$, puis que $A^{n}=I_2+n(A-I_2)$ pour tout
    $n\in\N$.
  \item En écrivant matriciellement la récurrence $u_{n+2}=2u_{n+1}-u_n$,
    retrouver que ses solutions sont les suites arithmétiques.
  \item Montrer que $S^{2}=-I_2$. En déduire que la suite de vecteurs
    $X_{n+1}=SX_n$ est périodique de période $4$, et interpréter la figure
    à l'aide de la question 3 : dans la base $(e_1,Se_1)$, $S$ agit comme
    le quart de tour $R$.
\end{questions}

\medskip
\begin{center}
\begin{tikzpicture}[scale=1.15,line cap=round,line join=round]
  \draw[otRule,line width=0.3pt] (-2.2,-1.6) grid (2.2,1.6);
  \draw[otMuted,line width=0.5pt,-{Stealth[length=1.6mm]}] (-2.3,0) -- (2.3,0);
  \draw[otMuted,line width=0.5pt,-{Stealth[length=1.6mm]}] (0,-1.7) -- (0,1.75);
  \fill[otGreen,opacity=0.12] (1,0) -- (1,1) -- (-1,0) -- (-1,-1) -- cycle;
  \draw[otGreenDark,line width=0.7pt,dash pattern=on 2.5pt off 1.5pt] (1,0) -- (1,1) -- (-1,0) -- (-1,-1) -- cycle;
  \foreach \p/\n/\a in {(1,0)/x/south east,(1,1)/Sx/south west,(-1,0)/S^{2}x/north west,(-1,-1)/S^{3}x/north east}{%
    \draw[otNavy,line width=1pt,-{Stealth[length=2mm]}] (0,0) -- \p;
    \node[font=\scriptsize,text=otNavy,anchor=\a] at \p {$\n$};}
\end{tikzpicture}\\[0.5mm]
\legende{L'orbite de $x=e_1$ sous $S$ : quatre vecteurs, sommets d'un
parallélogramme — l'image d'un carré, parcouru par un « quart de tour
déformé ».}
\end{center}
\end{defi}

